% ATTEMPT_STATUS: unresolved
% RESEARCH_GENERATED: 2026-10-01; GPT-6 Astra Ultra; individual mathematical attempt
% TOP_PROBLEM: {"id": 107, "title": "Corner Problem in Product Sets", "category_id": 2, "category": {"id": 2, "name": "combinatorics", "display_name": "Combinatorics", "description": "Counting problems, graph theory, discrete structures.", "slug": "combinatorics", "order_index": 2, "created_at": "2026-07-31T15:26:25.671Z"}, "status": "open"}
% FIRST_POSTED: 2026-10-01
% Imported from: unsolvedmath_additional_500.tex; UnsolvedMath ID 107
% !TeX program = lualatex
% Compile twice: lualatex 107.tex
% Self-contained: no external bibliography, images, or \input files.
% Numeric section labels and visible numbers are original UnsolvedMath IDs.
\documentclass[11pt,a4paper]{article}
\usepackage[margin=24mm,headheight=14pt]{geometry}
\usepackage{fontspec}
\setmainfont{Latin Modern Roman}
\setsansfont{Latin Modern Sans}
\setmonofont{Latin Modern Mono}
\usepackage{amsmath,amssymb,mathtools}
\usepackage{unicode-math}
\setmathfont{Latin Modern Math}
\usepackage{microtype}
\usepackage{enumitem,xurl,needspace}
\usepackage[dvipsnames]{xcolor}
\usepackage{hyperref}
\hypersetup{unicode=true,colorlinks=true,linkcolor=MidnightBlue,urlcolor=MidnightBlue,
 pdftitle={UnsolvedMath Problem 107},pdfauthor={Source-annotated compilation},bookmarksnumbered=true}
\usepackage{bookmark,tocloft,titlesec,fancyhdr}
\setlength{\cftsecnumwidth}{4.2em}
\cftsetpnumwidth{2.5em}
\cftsetrmarg{4em}
\titleformat{\section}{\Large\bfseries\raggedright}{\thesection}{0.7em}{}
\pagestyle{fancy}\fancyhf{}
\fancyhead[L]{\small\sffamily UnsolvedMath research attempt}
\fancyhead[R]{\small Original catalogue IDs}
\fancyfoot[C]{\thepage}
\setlength{\parindent}{0pt}
\setlength{\parskip}{5pt plus 1pt}
\setlength{\emergencystretch}{3em}
\setcounter{tocdepth}{1}\setcounter{secnumdepth}{1}
\setlist[itemize]{leftmargin=3em,itemsep=4pt,topsep=4pt,parsep=0pt}
\allowdisplaybreaks
\newcommand{\N}{\mathbb{N}}
\newcommand{\Z}{\mathbb{Z}}
\newcommand{\Q}{\mathbb{Q}}
\newcommand{\R}{\mathbb{R}}
\newcommand{\C}{\mathbb{C}}
\newcommand{\F}{\mathbb{F}}
\newcommand{\A}{\mathbb{A}}
\newcommand{\PP}{\mathbb{P}}
\newcommand{\E}{\mathbb{E}}
\newcommand{\eps}{\varepsilon}
\newcommand{\dd}{\,\mathrm d}
\newcommand{\sref}[2]{\hyperlink{source:#1:#2}{[S#2]}}
\newif\ifshowcatalogue
\showcataloguetrue
\newcommand{\cataloguescope}[1]{\ifshowcatalogue
 \paragraph{Statement recorded in the supplied source.}
 #1\fi}
\begin{document}
% UnsolvedMath ID 107; source catalogue code GREEN-018

\setcounter{section}{106}

\section{Corners in Products of Finite Groups}\label{107}

{\small\sffamily UnsolvedMath ID 107 · Combinatorics}

\subsection{Definitions and mathematical statement}

\paragraph{Shared notation.}Unless a section says otherwise, $\N=\{1,2,\ldots\}$,
$\Z$ is the ring of integers, $\Q,\R,\C$ are the rational, real, and complex
fields, $[N]=\{1,\ldots,\lfloor N\rfloor\}$, and $\log$ is the natural
logarithm. For real functions of a parameter tending to infinity,
$f\ll g$ and $f=O(g)$ mean $|f|\leq Cg$ eventually for some fixed $C$;
$f\gg g$ reverses this comparison; $f=o(g)$ means $f/g\to0$;
$f\sim g$ means $f/g\to1$; and $f\asymp g$ means both order bounds.
A subscript on $O$ or $\ll$ specifies allowed constant dependence.
The expression $n^{o(1)}$ in an upper bound means growth smaller than
$n^\epsilon$ for each fixed $\epsilon>0$. Local definitions govern any
exception, finite-field convention, or use of the same symbol for another
object. An asymptotic claim is not automatically an effective algorithm.



No commutativity of $G$ is assumed, and $gx,gy$ denote left multiplication. For $A\subseteq G^2$ put $\alpha=|A|/|G|^2$ and
\[C(A)=\sum_{x,y,g\in G}1_A(x,y)1_A(gx,y)1_A(x,gy).\]
The notation $C(A)\gg_\alpha |G|^3$ asks for a positive constant depending only on $\alpha$, uniformly over the groups. The source includes $g=1$; those degenerate triples number at most $|G|^2$. \sref{107}{1} \sref{107}{2}

\paragraph{Statement recorded in the supplied source.}
Suppose $G$ is a finite group, and let $A \subset G \times G$ be a subset of density $\alpha$. Are there $\gg_\alpha |G|^3$ triples $x, y, g$ such that $(x, y), (gx, y), (x, gy)$ all lie in $A$? \sref{107}{1}

\paragraph{Residual task reported by the archive.}

The embedded review (2026-08-17) records: Prove or refute a uniform c(alpha)|G|\textasciicircum{}3 lower bound for nontrivial naïve corners in every sufficiently large finite group. \sref{107}{1}

\subsection{Short English statement}

Does density alone force many multiplicative corner configurations in the Cartesian square of every finite group?

\subsection{Research attempt}

\paragraph{Provenance and scope.}
This individual attempt was generated by GPT-6 Astra Ultra on 1 October 2026. The method was to derive explicit partial results and test the first proposed extension adversarially. The original statement and sources below are retained. No claim of mathematical novelty or complete solution is made.

\paragraph{Formal target and source distinction.}
Let $G$ be any finite group of order $n$, with multiplication written in the displayed order, and let $A\subset G^2$ have density $\alpha$. Write $C(A)$ for the number of triples $(x,y,g)$ for which $(x,y),(gx,y),(x,gy)\in A$. The primary source distinguishes these naive corners from configurations using right multiplication in one coordinate. That distinction is retained throughout: commutativity is never assumed. We seek a lower bound depending only on density and prove one above a concrete threshold, together with a separate result for Cartesian products.

\paragraph{A row-column lower bound.}
Let $R_y=\{x:(x,y)\in A\}$ and $S_x=\{y:(x,y)\in A\}$, with sizes $r_y,s_x$. For a fixed occupied point $(x,y)$, the possible multipliers for the second corner point form $R_yx^{-1}$, and those for the third point form $S_xy^{-1}$. These are subsets of the same group, of sizes $r_y$ and $s_x$. Their intersection has size at least $r_y+s_x-n$. Summing over occupied points yields
\[
 C(A)\geq\sum_y r_y^2+\sum_xs_x^2-n|A|.
\]
The right-hand side is allowed to be negative; the true count is nonnegative. Cauchy--Schwarz gives $\sum_y r_y^2\geq |A|^2/n$ and the analogous column inequality. Thus
\[
 C(A)\geq\alpha(2\alpha-1)n^3.
\]
For every fixed $\alpha>1/2$ this answers the counting question affirmatively in that density range, by a completely elementary argument. If nonidentity multipliers are required, remove exactly $|A|=\alpha n^2$ triples, since $g=1$ always works for each occupied point. Hence for sufficiently large $n$ the nontrivial count is at least $\tfrac12\alpha(2\alpha-1)n^3$.

\paragraph{A more informative form retains degree variation.}
Put $\bar r=\bar s=\alpha n$. Expanding the squares proves
\[
 C(A)\geq\alpha(2\alpha-1)n^3
 +\sum_y(r_y-\alpha n)^2+\sum_x(s_x-\alpha n)^2.
\]
Thus degree irregularity helps this method rather than hurting it. In particular it proves a positive-density lower bound below one half whenever the two variance sums exceed $\alpha(1-2\alpha)n^3$ by a fixed positive multiple of $n^3$. The unhandled sets for this argument are not merely sparse: they must also have fairly balanced row and column sizes. This is a precise restriction on any difficult configuration, albeit far from a complete structural description.

\paragraph{Cartesian products admit a different bound at every density.}
Suppose $A=B\times D$, with $B,D\subset G$. A corner is equivalent to $b_1,b_2\in B$, $d_1,d_2\in D$, with
\[
 b_2b_1^{-1}=d_2d_1^{-1}.
\]
Rearranging in the correct order gives $d_2^{-1}b_2=d_1^{-1}b_1$. Let $r(t)$ count pairs $(d,b)\in D\times B$ with $d^{-1}b=t$. Every equality of two such pairs determines a unique multiplier $g=b_2b_1^{-1}$ and a unique corner. Therefore
\[
 C(B\times D)=\sum_{t\in G}r(t)^2
 \geq\frac{(|B||D|)^2}{n}=\alpha^2n^3.
\]
The inequality is Cauchy--Schwarz and uses only $\sum_t r(t)=|B||D|$. This proof works for highly noncommutative groups. Replacing $d^{-1}b$ by $bd^{-1}$ without redoing the calculation would be an ordering error.

\paragraph{Why these two results do not solve the low-density case.}
A general dense bipartite set need not contain a Cartesian rectangle of positive relative sizes in both coordinates uniformly as $n$ grows. Finite small rectangles do not help: a rectangle whose area is $o(n^2)$ gives only an $o(n^3)$ bound through the displayed formula. Decomposing $A$ into rectangles is also not automatically useful, because the lower bounds are quadratic in their densities and their total can vanish when the number of pieces grows. The row-column argument fails for regular sparse examples because its lower bound becomes nonpositive. Combining the two observations thus requires a genuine structural or correlation theorem, which has not been established here.

\paragraph{Verification and status.}
For $A=G^2$, both arguments give $n^3$, the exact count. For a single point, the true count is one and the row-column lower estimate is merely weak, not contradictory. For $B=D=H$ a subgroup, the representation formula gives $|H|^3$, again the exact answer. These boundary checks test the multiplication order and the identity contribution. The problem remains unresolved for arbitrary fixed densities at most one half. The checked results are the uniform high-density theorem, its variance refinement, and the all-density Cartesian-product theorem.

\subsection{Sources}

{\small

\begin{itemize}

\item[\hypertarget{source:107:1}{S1}] ULAM.AI, \emph{UnsolvedMath}, supplied \texttt{problems.json}, record 107 (GREEN-018), statement and background. The embedded literature-review date is 2026-08-17; that is the archive’s review date, not a fresh certification. Dataset landing page: \url{https://huggingface.co/datasets/ulamai/UnsolvedMath}.

\item[\hypertarget{source:107:2}{S2}] Tim Austin, Ajtai-Szemerédi Theorems over quasirandom groups, arXiv:1503.08746; Recent Trends in Combinatorics (2016), 453-484. \url{https://arxiv.org/abs/1503.08746}. Bibliographic information imported from the supplied record.

\item[\hypertarget{source:107:3}{S3}] Michael Jaber, Yang P. Liu, Shachar Lovett, Anthony Ostuni, and Mehtaab Sawhney, Quasipolynomial bounds for the corners theorem, arXiv:2504.07006. \url{https://arxiv.org/abs/2504.07006}. Bibliographic information imported from the supplied record.

\item[\hypertarget{source:107:4}{S4}] Ben Green, 100 Open Problems, Problem 18, . \url{https://people.maths.ox.ac.uk/greenbj/papers/open-problems.pdf}. The author-hosted document was inspected on 27 September 2026; the archive’s local code is not a problem-number locator in this paper.

\item[\hypertarget{source:107:5}{S5}] Formal Conjectures, Ben Green Open Problem 18. \url{https://firsching.ch/formal-conjectures/src/FormalConjectures/GreensOpenProblems/%C2%AB18%C2%BB/}. Bibliographic information imported from the supplied record.

\end{itemize}

}

\label{end:107}
\end{document}
